\documentclass[options]{philosophersimprint}
%\usepackage{sectsty}
%\linespread{1.2}
\usepackage{xunicode}
\usepackage{xltxtra}
\defaultfontfeatures{Mapping=tex-text}
\setmainfont[
Ligatures={Common}, Numbers={OldStyle}, Variant=01,
BoldFont=LinLibertine_RB.otf,
ItalicFont=LinLibertine_RI.otf,
BoldItalicFont=LinLibertine_RBI.otf
]{LinLibertine_R.otf}
\setmonofont[Scale=0.8]{DejaVuSansMono.ttf}
\usepackage{stmaryrd}
\usepackage[toc,page]{appendix}
\usepackage{amsmath,amsthm,amssymb}
\usepackage[colorlinks=true, citecolor=blue, linkcolor=blue, hyperfootnotes=false]{hyperref}
\usepackage{nameref}
\usepackage{soul}
\usepackage{frame}
\usepackage[style=apa, natbib=true]{biblatex}
\usepackage{scrextend}
\usepackage{mathrsfs}
%\addtokomafont{labelinglabel}{\sffamily}
\makeatletter
\newcommand\labelitem[2]{%
  \item[#1]\def\@currentlabel{#1}\label{#2}}
\makeatother
\usepackage[inline]{enumitem}
\usepackage{framed, csquotes}
\usepackage{tikz}
\usepackage{scalerel}
\usepackage{titlesec}
\newcommand\halo{{{\circ}\mkern-1mu}}
\usepackage{linguex}
\usepackage{parskip}
\usepackage{fancyhdr}
\pagestyle{plain}
% Diagrams
%\usepackage{pgfplots}
\usepackage{tikz}
\usepackage{qtree}

\usepackage{fitch}



\usepackage{FOL_Yale/notation}

\begin{document}

PHIL 1115: First-order Logic \hfill Yale, Fall 2026\\
Lecture 10, Handout \hfill Dr. Daniel Grimmer

\textbf{1. Inferentialism, Revisited} Recall last time's idea: to know what a connective \textit{means} is to know how to correctly \textit{use} it in argument: which inferences into it, and which inferences out of it, are licensed. For ``and'', ``inclusive-or'' and ``material biconditional'' (formalized as $\ConjTight$, $\Disj$, and $\Bicond$) this was easy: (almost) everybody agrees. But the situation is different for ``if ... then'' and ``it is not the case that ...'' (formalized as $\Cond$ and $\Neg$). What does it really take to correctly infer that $A \Cond B$? And what does it really take to correctly infer $A$ from $\Neg\Neg A$? These turn out to be genuinely contested, as we shall see.

My whole job today is to teach you the standard, classical rules for $\Cond$ and $\Neg$. \textbf{These standard rules are the ones you need to know for the mid-term and final}, so we'll get them all cleanly onto the table. As we go, I'll flag (but not yet pursue) exactly where each of these rules becomes controversial: there are two famous forks in the road ahead. Those forks are the whole business of \textit{next} lecture (and, to be clear, none of that critique is on the exam).

\begin{center}
\begin{minipage}[t]{0.26\textwidth}
\textbf{2. The Conditional}
\vspace{-0.25cm}
\begin{align*}
\begin{nd}
\hypo{1}{p \Disj q}
\hypo{2}{p \Cond r}
\hypo{3}{q \Cond s}
\open
\hypo{4}{p}
\have{5}{r\qquad\CondE,2,4}
\have{6}{r \Disj s\quad\DisjI,5}
\close
\open
\hypo{7}{q}
\have{8}{s\qquad\CondE,3,7}
\have{9}{r \Disj s\quad\DisjI,8}
\close
\have{10}{r \Disj s\qquad \DisjE,1,4-6,7-9}
\end{nd}
\end{align*}
\end{minipage}
\hfill
\begin{minipage}[t]{0.19\textwidth}
\vspace{0.5cm}
Here is an English argument for what is often called the \textbf{constructive dilemma}: Our premises are that $p \Disj q$ and $p \Cond r$ and $q \Cond s$. We would like to conclude that $r \Disj s$.
\vspace{0.25cm}\\
In the $p$ case, modus ponens gets us $r$ which we can weaken to $r\Disj s$. Similarly, in the $q$ case, modus ponens gets us $s$ which we can weaken to $r\Disj s$. Hence, either way we get $r\Disj s$.
\end{minipage}
\end{center}

Lines 5 and 8 of this argument both use the elimination rule for $\Cond$. You already know it under a different name: \textbf{modus ponens}. Nobody disputes this rule.

\begin{center}
\begin{minipage}[t]{0.20\textwidth}
\textbf{$\Cond$-Elimination Rule:}
\vspace{-0.25cm}
\begin{align*}
 \begin{nd}
        \have[n]{}{A \Cond B}
        \have[...]{}{\dots}
        \have[m]{}{A}
        \have[...]{}{\dots}
        \have[k]{}{B    \qquad\CondE,n,m}
    \end{nd}
\end{align*}
\end{minipage}
\hfill
\begin{minipage}[t]{0.20\textwidth}
\textbf{$\Cond$-Introduction Rule:}
\vspace{-0.25cm}
\begin{align*}
\begin{nd}
\open
\hypo[n]{1}{A}
\have[...]{}{\dots}
\have[m]{}{B}
\close
\have[m+1]{}{A \Cond B    \qquad\CondI,n,m}
\end{nd}
\end{align*}
\end{minipage}
\end{center}
But what does it take to \textit{derive} $A \Cond B$ in the first place? The classical answer is given above: open a sub-derivation, assume $A$, and see whether you can derive $B$ inside it. If you can, you're entitled to close the sub-derivation, discharge the assumption, and conclude $A \Cond B$ outright.

Here's that rule earning its keep by validating the \textbf{hypothetical syllogism}.
\begin{center}
\vspace{-0.75cm}
\begin{minipage}[t]{0.25\textwidth}
\begin{align*}
\begin{nd}
\hypo{1}{p \Cond q}
\hypo{2}{q \Cond r}
\open
\hypo{3}{p}
\have{4}{q\qquad\CondE,1,3}
\have{5}{r\qquad\CondE,2,4}
\close
\have{6}{p \Cond r\qquad\CondI,3,5}
\end{nd}
\end{align*}
\end{minipage}
\hfill
\begin{minipage}[t]{0.20\textwidth}
\vspace{0.5cm}
Here is the corresponding English argument: Our premises are that $p \Cond q$ and $q \Cond r$. We would like to conclude that $p \Cond r$.
\vspace{0.25cm}\\
Suppose $p$. Then by applying modus ponens twice we get $q$ and then $r$. Hence, if we assume $p$, then we can conclude $r$. That is to say, $p\Cond r$.
\end{minipage}
\end{center}
\vspace{-0.5cm}
\textbf{So far, so good.} But keep one eye on $\Cond$-Introduction as we go. In both proofs above, the antecedent we assumed was genuinely \textit{used} on the way to the consequent (each time, we fed it into modus ponens). Ask yourself: does the rule, as I stated it, actually \textit{require} that? Let's see an example where the antecedent is completely un-used.

Before this, however, we need one small, almost trivial rule that we've somehow managed to avoid until now.
\begin{center}
\begin{minipage}[t]{0.25\textwidth}
\textbf{The Reiteration Rule:}
\vspace{-0.5cm}
\begin{align*}
    \begin{nd}
        \have[n]{} {A}
        \have[...]{} {\dots}
        \have[m]{} {A\qquad\qquad \Reit,n}
    \end{nd}
\end{align*}
\end{minipage}
\hfill
\begin{minipage}[t]{0.20\textwidth}
Informally: if you have a formula $A$ on an accessible line, you're always entitled to simply copy it down again on any later accessible line (in particular, down into a sub-derivation).
\end{minipage}
\end{center}
\vspace{-0.5cm}
That's the whole rule. We'll use it immediately to prove, that $p\ProvesND q \Cond p$.
\vspace{-0.75cm}
\begin{center}
\begin{minipage}[t]{0.25\textwidth}
\begin{align*}
\begin{nd}
\hypo{1}{p}
\open
\hypo{2}{q}
\have{3}{p\qquad\Reit,1}
\close
\have{4}{q \Cond p\qquad\CondI,2,3}
\end{nd}
\end{align*}
\end{minipage}
\hfill
\begin{minipage}[t]{0.20\textwidth}
\vspace{0.5cm}
Here is the corresponding English argument: Given our premise of $p$, assume $q$. Then recall that $p$. Hence, we assumed $q$ and got to $p$. That is to say, $q\Cond p$ as desired.
\end{minipage}
\end{center}
\vspace{-0.25cm}
This is one of the ``paradoxes of the material conditional'' which we discussed in Lecture 3. Plugging in a true sentence for $p$ yields a logical tautology \textit{no matter what we pick for} $q$. Hence, the truth of the following statements is logically necessary:
\begin{itemize}
    \item If $1+1=2$ then the fact that ``the sky is blue'' implies that $1+1=2$.
    \item If $2+2=4$ then the fact that ``I am the Pope'' implies that $2+2=4$.
\end{itemize}

To avoid this kind of strangeness, some logicians have become suspicious of this kind of argumentation. Notice that to derive $q \Cond p$, we opened a sub-derivation, assumed $q$, and then \textit{never used it}: line 3 simply reiterated $p$ from outside, and $q$ sat there doing nothing. We've proven ``if $q$, then $p$'' with $q$ having contributed nothing whatsoever. While \textit{this maneuver is classically allowed}, some logicians take issue with it (more on them next lecture).

\textbf{3. Negation} The rules for $\Neg$ come in a pair too (introduction and elimination) along with one piece of new notation Over and over in what follows we will reach some formula $B$ together with its negation $\Neg B$, and want to say: \textit{something has gone horribly wrong here}. It does not matter what particular contradiction we have hit (e.g., $B$ and $\Neg B$) just that it is some contradiction or other. \textit{All contradictions are alike} and we use one symbol to mark them all: the \textbf{falsum}, written $\Falsum$. Here is its introduction rule.
\begin{center}
\begin{minipage}[t]{0.2\textwidth}
\textbf{$\Falsum$-Introduction Rule:}
\vspace{-0.25cm}
\begin{align*}
\begin{nd}
\have[n]{}{B}
\have[...]{}{\dots}
\have[m]{}{\Neg B}
\have[...]{}{\dots}
\have[k]{}{\Falsum\qquad\quad\FalsumI,n,m}
\end{nd}
\end{align*}
\end{minipage}
\hfill
\begin{minipage}[t]{0.2\textwidth}
\textbf{$\Neg$-Introduction Rule:}
\vspace{-0.25cm}
\begin{align*}
\begin{nd}
\open
\hypo[n]{1}{A}
\have[...]{}{\dots}
\have[m]{2}{\Falsum}
\close
\have{}{\Neg A\qquad\quad\NegI,n,m}
\end{nd}
\end{align*}
\end{minipage}
\end{center}
\vspace{-0.5cm}
Note that $\Falsum$ is not a part of our language, you will never see it \textit{inside} a formula. Concretely, $p\Disj\Falsum$ and $\Falsum\Bicond(q\Conj\Neg q)$ are not formulas of our language (not even unofficially). You will only ever see $\Falsum$ in the context of natural derivations, standing alone on its own line to indicate that a certain line of reasoning has hit a dead end.

This symbol is used in the $\Neg$-introduction rule. To derive $\Neg A$, assume $A$, and argue your way to $\Falsum$. \textbf{Nobody disputes this rule}: if accepting $A$ leads straight to contradiction, everybody agrees $A$ should be rejected. Next lecture we will meet logicians who go further and treat this rule as the very \textit{definition} of negation, which tells you how uncontroversial it is.

To see this rule in action, recall this argument from Kant from Lecture 4:
\begin{quote}
If a moral theory is studied empirically then examples of conduct will be considered. And if examples of conduct are considered, principles for selecting examples will be used. But if principles for selecting examples are used, then moral theory is not being studied empirically. Therefore, moral theory is not studied empirically.
\end{quote}
\vspace{-0.5cm}
$$p\Cond q, \quad q\Cond r, \quad r\Cond \Neg p \quad \therefore \quad \Neg p$$
\begin{center}
\begin{minipage}[t]{0.27\textwidth}
\vspace{-1.25cm}
\begin{align*}
\begin{nd}
\hypo{1}{p \Cond q}
\hypo{2}{q \Cond r}
\hypo{3}{r \Cond \Neg p}
\open
\hypo{4}{p}
\have{5}{q\qquad\quad\CondE,1,4}
\have{6}{r\qquad\quad\CondE,2,5}
\have{7}{\Neg p\qquad\CondE,3,6}
\have{8}{p\qquad\quad\Reit,4}
\have{9}{\Falsum\qquad\ \FalsumI,8,7}
\close
\have{10}{\Neg p\qquad\ \NegI,4,9}
\end{nd}
\end{align*}
\end{minipage}
\hfill
\begin{minipage}[t]{0.18\textwidth}
Here is the corresponding English argument. Our premises are $p \Cond q$ and $q\Cond r$ and $r\Cond \Neg p$. We hope to conclude $\Neg p$.
\vspace{0.25cm}\\
To start we assume $p$. But three applications of modus ponens then yield $q$ and $r$ and $\Neg p$. Recalling that we began from $p$, this is a contradiction, $\Falsum$. Hence we must reject our assumption of $p$ (adding a $\Neg$).
\end{minipage}
\end{center}
\vspace{-0.25cm}
\begin{center}
\begin{minipage}[t]{0.27\textwidth}
Notice that the $\Neg$-introduction rule only ever \textit{adds} negations to our formulas. We need a rule for taking them off again: \textit{negations come off in pairs}. (Hence the term \textbf{double negation elimination}, or \textbf{DNE}.)
\end{minipage}
\hfill
\begin{minipage}[t]{0.15\textwidth}
\textbf{$\Neg$-Elimination Rule:}
\vspace{-0.25cm}
\begin{align*}
\begin{nd}
\have[n]{}{\Neg\Neg A}
\have[...]{}{\dots}
\have[m]{}{A\qquad\NegE,n}
\end{nd}
\end{align*}
\end{minipage}
\end{center}
\vspace{-0.25cm}
Perhaps surprisingly, this is one of the most contested rules in modern logic. We will see why next lecture. Here is a simple use case:
\begin{center}
\vspace{-0.75cm}
\begin{minipage}[t]{0.2\textwidth}
\begin{align*}
\begin{nd}
\hypo{1}{\Neg p \Cond q}
\hypo{2}{\Neg p \Cond \Neg q}
\open
\hypo{3}{\Neg p}
\have{4}{q\qquad\ \ \CondE,1,3}
\have{5}{\Neg q\qquad\CondE,2,3}
\have{6}{\Falsum\qquad\ \FalsumI,4,5}
\close
\have{7}{\Neg\Neg p\quad\ \ \NegI,3,6}
\have{8}{p\qquad\qquad\NegE,7}
\end{nd}
\end{align*}
\end{minipage}
\hfill
\begin{minipage}[t]{0.23\textwidth}
\vspace{0.5cm}
Here is the corresponding English argument. Our premises are $\Neg p \Cond q$ and $\Neg p \Cond \Neg q$ and we hope to conclude $p$. To start we assume $\Neg p$. But two applications of modus ponens then yield $q$ and $\Neg q$, which is a contradiction, $\Falsum$. Hence we must reject our assumption of $\Neg p$, giving $\Neg\Neg p$. Then we strip the two negations off to get $p$.
\end{minipage}
\end{center}

\begin{center}
\begin{minipage}[t]{0.21\textwidth}
These last three steps will occur together very frequently. In fact, this pattern is so common that it has a name, \textbf{classical reductio}, and it is worth writing down once in general form: If $\Neg A$ leads to a contradiction, then $\Neg\Neg A$ and hence $A$ by double negation elimination.
\end{minipage}
\hfill
\begin{minipage}[t]{0.21\textwidth}
\vspace{-0.5cm}
\begin{align*}
\begin{nd}
\open
\hypo[n]{1}{\Neg A}
\have[...]{}{\dots}
\have[m]{2}{\Falsum}
\close
\have{3}{\Neg\Neg A\qquad\NegI,n,m}
\have{}{A\qquad\quad \NegE,m+1}
\end{nd}
\end{align*}
\end{minipage}
\end{center}
\vspace{-0.25cm}
Importantly, this is \textit{not} a basic rule in our ND system. If you ever want to remove a single negation, you must find a way to add a second negation and then remove the pair.

Using $\Neg$-elimination we can derive the \textbf{Law of the Excluded Middle (LEM)} from no premises: $\ProvesND p\Disj\Neg p$.
\begin{center}
\vspace{-0.75cm}
\begin{minipage}[t]{0.2\textwidth}
\begin{align*}
\begin{nd}
\open
\hypo{1}{\Neg(p \Disj \Neg p)}
\open
\hypo{2}{p}
\have{3}{p \Disj \Neg p\qquad\quad\ \DisjI,2}
\have{4}{\Neg(p \Disj \Neg p)\qquad\ \Reit,1}
\have{5}{\Falsum\qquad\quad\quad\ \FalsumI,3,4}
\close
\have{6}{\Neg p\qquad\qquad\quad\ \ \NegI,2,5}
\have{7}{p \Disj \Neg p\qquad\qquad\ \ \DisjI,6}
\have{8}{\Neg(p \Disj \Neg p)\qquad\ \ \Reit,1}
\have{9}{\Falsum\qquad\qquad\quad\ \ \FalsumI,7,8}
\close
\have{10}{\Neg\Neg(p \Disj \Neg p)\ \qquad \NegI,1,9}
\have{11}{p \Disj \Neg p\qquad\qquad\NegE,10}
\end{nd}
\end{align*}
\end{minipage}
\hfill
\begin{minipage}[t]{0.18\textwidth}
\vspace{0.5cm}
Here is the corresponding English argument: Assume LEM fails, $\Neg(p \Disj \Neg p)$. Then assume $p$. But this can be weakened to $p \Disj \Neg p$, which contradicts our initial assumption. Because our assumption of $p$ led to a contradiction, we have $\Neg p$. But this too can be weakened to $p \Disj \Neg p$, yielding another contradiction. Hence we must reject our initial assumption, giving $\Neg\Neg(p \Disj \Neg p)$ at line 10, and strip its two negations at line 11.
\end{minipage}
\end{center}
\vspace{-0.25cm}

We now have \textbf{twelve} rules in all. Ten of them come in pairs (introduction and elimination) and the other two are for bookkeeping: Reiteration for where a formula may be reused, and $\Falsum$I for when a contradiction has been reached.
%\vspace{0.25cm}
\begin{framed}
\noindent\textbf{Basic Rule Summary.} All twelve, one line each (see the Natural Deduction Cheat Sheet for the full Fitch-style schemas):\\[-8pt]
\begin{tabular}{@{}rl@{}}
$\Reit$ & from $A$: derive $A$ (on any accessible line)\\
$\ConjI$ & from $A$ and $B$: derive $A\Conj B$\\
$\ConjE$ & from $A\Conj B$: derive $A$, or derive $B$\\
$\DisjI$ & from $A$ (or from $B$): derive $A\Disj B$\\
$\DisjE$ & from $A\Disj B$, and two proofs of $C$ (from both $A$ and $B$): derive $C$\\
$\BicondI$ & from a proof of $B$ from $A$, and a proof of $A$ from $B$: derive $A\Bicond B$\\
$\BicondE$ & from $A\Bicond B$ and $A$ (or $B$): derive $B$ (or $A$)\\
$\CondI$ & from a proof of $B$ from the assumption $A$: derive $A\Cond B$\\
$\CondE$ & from $A\Cond B$ and $A$: derive $B$\\
$\NegI$ & from a proof of $\Falsum$ from the assumption $A$: derive $\Neg A$\\
$\NegE$ & from $\Neg\Neg A$: derive $A$ \quad(double negation elimination)\\
$\FalsumI$ & from $B$ and $\Neg B$: derive $\Falsum$\\
\end{tabular}
\vspace{-0.75cm}
\end{framed}

\textbf{4. Completing the ND Dictionary}\\
Recall the definitions we already have from Lecture 9: ND-valid (denoted $X \ProvesND A$), ND-invalid (denoted $X \not\ProvesND A$), ND-tautology (denoted $\ProvesND A$), and ND-equivalent. Aligning our new method of natural deduction with the table and tree methods yields several more definitions:

\textbf{Contradiction:} A set of formulas $X$ is jointly an \textbf{ND-contradiction} (denoted $X \ProvesND$) if and only if $X \ProvesND \Falsum$. That is, if one can derive $\Falsum$ from $X$. For a single formula, $A$, this is equivalent to saying that $\Neg A$ is an ND-tautology. Namely, $A \ProvesND$ if and only if $\ProvesND\Neg A$. (Why? See the following proofs.)
\begin{center}
\begin{minipage}[t]{0.10\textwidth}
\vspace{-0.75cm}
\begin{align*}
\begin{nd}
\open
\hypo[1]{1}{A}
\have[...]{}{\dots}
\have[n]{2}{\Falsum}
\close
\have{3}{\Neg A\qquad\NegI,1,n}
\end{nd}
\end{align*}
\end{minipage}
\hfill
\begin{minipage}[t]{0.1\textwidth}
\vspace{-0.75cm}
\begin{align*}
\begin{nd}
    \hypo{1}{A}
    \have[...]{}{\dots \quad \text{ Proof of } \Neg A}
    \have[k]{}{\Neg A}
    \have[k+1]{}{A \qquad R, 1}
    \have[k+2]{}{\Falsum \qquad \FalsumI,k+1,k}
\end{nd}
\end{align*}
\end{minipage}
\end{center}

Left) Suppose that one can derive $\Falsum$ from $A$. Then begin a proof with no premises, assume $A$, derive $\Falsum$, and then use $\Neg$-Introduction to conclude $\Neg A$. This demonstrates that $\Neg A$ is an ND-tautology.

Right) Suppose one has a zero-premise proof of $\Neg A$. Begin an argument with premise $A$, and derive $\Neg A$ from no assumptions. But this contradicts our premise of $A$, hence $\Falsum$. This demonstrates that $A$ is an ND-contradiction.

\textbf{Satisfiable}: A set of formulas $X$ is jointly \textbf{ND-satisfiable} (denoted $X \not\ProvesND$) if and only if they are \textit{not} jointly an ND-contradiction: that is, there is no way whatsoever to derive $\Falsum$ from $X$.

\textbf{Falsifiable}: A formula $A$ is \textbf{ND-falsifiable} (denoted $\not\ProvesND A$) if and only if it is \textit{not} an ND-tautology: that is, there is no way whatsoever to derive $A$ from no assumptions.

\textbf{Contingent}: A formula $A$ is \textbf{ND-contingent} if and only if it is both ND-satisfiable and ND-falsifiable.

As with ND-invalidity, there is no \textit{direct} way to demonstrate these last three properties: They all require showing that a certain kind of derivation is not possible. But, in practice, the space of all possible derivations is far too unwieldy to search by hand. The only way to test for these ND properties is indirectly via the table and tree methods (more on this in Lecture 12).

\begin{framed}
\noindent\textbf{Definitions Summary: ND-Notions and Notation}\\[-12pt]
% @{} on both ends and the trimmed \tabcolsep are load-bearing: at the default
% 6pt with ordinary outer padding this tabular runs 14pt wider than the column,
% and the \hline between the two halves pokes out past the frame's right rule.
\setlength{\tabcolsep}{5pt}
\begin{tabular}{@{}l | l@{}}
Direct Confirmation Possible:& \\
ND-valid ($X\ProvesND A$) & Derive $A$ from $X$\\
ND-tautology ($\ProvesND A$) & Derive $A$ from nothing\\
ND-contradiction ($X\ProvesND$) & Derive $\Falsum$ from $X$\\
ND-equivalent & Derive $B$ from $A$ and $A$ from $B$\\
\hline
Only Indirect Confirmation:& \\
ND-invalid ($X\not\ProvesND A$) & No derivation of $A$ from $X$ exists\\
ND-satisfiable ($X\not\ProvesND$) & No derivation of $\Falsum$ from $X$ exists\\
ND-falsifiable ($\not\ProvesND A$) & No derivation of $A$ from nothing\\
ND-contingent & Both ND-satisfiable and ND-falsifiable\\
\end{tabular}
% -0.25cm is the most that can be clawed back here: at -0.35 and beyond
% the frame's bottom rule rides up through the ND-contingent row and leaves it
% outside the box. Anything smaller pushes the handout to five pages.
\vspace{-0.25cm}
\end{framed}

\end{document}
